\subsection{一元多项式的根。重数}

设 \(g \in A[(X_i)_{i \in I}]\) ，并设 \(E\) 是一个含单位元的结合 \(A\) -代数。设 \(x = (x_i)_{i \in I}\) 是 \(E\) 中两两可交换的元素族。我们就说 \(x\) 是 \(g\) 在 \(E^{I}\) 中的零点，如果 \(g(x) = 0\) 。若 \(f\) 是单个未定元的多项式，则 \(f\) 在 \(E\) 中的零点也称为 \(f\) 在 \(E\) 中的根。

命题1。——设 \(f \in \mathbf{A}[\mathbf{X}]\) 与 \(\alpha \in \mathbf{A}\) 。 \(f\) 除以 \(\mathbf{X} - \alpha\) 的余式是 \(f(\alpha)\) 。 \(\alpha\) 成为 \(f\) 的根的充分必要条件是 \(\mathbf{X} - \alpha\) 是 \(f\) 在 \(\mathbf{A}[\mathbf{X}]\) 中的因子。

因为如果 \(u, v \in \mathbf{A}[\mathbf{X}]\) 使得 \(f = (\mathbf{X} - \alpha)u + v\) , \(\deg v < 1\) ，则 \(v\) 是一个标量且 \(f(\alpha) = (\alpha - \alpha)u(\alpha) + v = v\) 。这证明了第一个断言，第二个断言由它推出。

命题2。——设 \(f \in \mathbf{A}[\mathbf{X}]\) ， \(\alpha \in \mathbf{A}\) ，并设 \(h\) 为整数 \(\geqslant 0\) 。以下条件等价：

\begin{enumerate}
\def\labelenumi{(\roman{enumi})}
\item
  \(f\) 能被 \((\mathbf{X} - \alpha)^h\) 整除，但不能被 \((\mathbf{X} - \alpha)^{h + 1}\) 整除；
\item
  存在 \(g \in \mathbf{A}[\mathbf{X}]\) 使得 \(f = (\mathbf{X} - \alpha)^h g\) 且 \(g(\alpha) \neq 0\) 。
\item
  \(\Rightarrow\) (ii) 由命题 1 立即可得。
\item
  \(\Rightarrow\) (i)：假设 \(f = (\mathbf{X} - \alpha)^h g\) ，其中 \(g\) 不以 \(\alpha\) 为根。那么 \(f\) 能被 \((\mathbf{X} - \alpha)^h\) 整除；如果存在 \(g_1 \in \mathbf{A}[\mathbf{X}]\) 使得 \(f = (\mathbf{X} - \alpha)^{h + 1}g_1\) ，则由于 \((\mathbf{X} - \alpha)^n\) 在 \(\mathbf{A}[\mathbf{X}]\) （IV，第 9 页，命题 7）中不是零因子，我们有 \(g = (\mathbf{X} - \alpha)g_1\) ，从而 \(g(\alpha) = 0\) ，这是荒谬的。
\end{enumerate}

命题3. --- 设 \(f\) 是 \(\mathbf{A}[\mathbf{X}]\) 的非零元素，且 \(\alpha \in \mathbf{A}\) 。存在唯一一个整数 \(h \geqslant 0\) 满足命题 2 的条件 (i) 与 (ii) 。

这一点对条件 (i) 是显然的，只要注意到：若 \(f\) 能被 \((\mathbf{X} - \alpha)^h\) 整除，则 \(\deg f \geqslant h\) （IV，第 9 页，命题 7）。

定义1. --- 在上述记号下，我们说 \(\alpha\) 的阶为 \(h\) ，或重数为 \(h\) （相对于 \(f\) ）。

若 h \textgreater{} 0 ，我们也称 \(\alpha\) 为 f 的 h 阶根或重数 h 的根。1 阶根称为单根，2 阶根称为二重根，\ldots{} 阶数 \textgreater{} 1 的根称为重根。

注记。--- 1) 若 \(f=0\) ，我们约定说 \(\alpha\) 的阶 \(\geqslant h\) （相对于 \(f\) ），不论 \(\alpha\in\mathbf{A}\) 与整数 \(h\geqslant0\) 如何。对任意 \(f\in\mathbf{A}[\mathbf{X}]\) 与 \(\alpha\in\mathbf{A}\) ，说 \(\alpha\) 的阶 \(\geqslant h\) （相对于 \(f\) ）意味着 \((\mathbf{X}-\alpha)^{h}\) 整除 \(f\) 。

\begin{enumerate}
\def\labelenumi{\arabic{enumi})}
\setcounter{enumi}{1}
\tightlist
\item
  设 B 为包含 A 作为子环的交换环。设 \(f \in A[X]\) 非零且 \(\alpha \in A\) 。 \(\alpha\) 相对于 f 的阶是相同的（无论把 f 视为 B{[}X{]} 的元素还是 A{[}X{]} 的元素）。这由命题 2 的条件 (ii) 可以清楚地看出。
\end{enumerate}

命题 4. --- 设 f 和 g 是 A{[}X{]} 的非零元素。设 \(\alpha \in \mathbf{A}\) ，并设 \(\alpha\) 相对于 f 和 g 的阶分别为 p 和 q。

\begin{enumerate}
\def\labelenumi{(\roman{enumi})}
\item
  \(\alpha\) 相对于 \(f + g\) 的阶 \(\geqslant \inf(p, q)\) ，且它等于 \(\inf(p, q)\) （当 \(p \neq q\) 时）。
\item
  \(\alpha\) 相对于 \(fg\) 的阶 \(\geqslant p + q\) ，且它等于 \(p + q\) （当 \(\mathbf{A}\) 是整环时）。
\end{enumerate}

因为我们有 \(f(\mathbf{X}) = (\mathbf{X} - \alpha)^p f_1(\mathbf{X}), g(\mathbf{X}) = (\mathbf{X} - \alpha)^q g_1(\mathbf{X})\) ，其中 \(f_1(\alpha) \neq 0\) , \(g_1(\alpha) \neq 0\) 。例如设 \(p \leqslant q\) ；于是我们有

\[
f (\mathbf {X}) + g (\mathbf {X}) = (\mathbf {X} - \alpha) ^ {p} \left(f _ {1} (\mathbf {X}) + (\mathbf {X} - \alpha) ^ {q - p} g _ {1} (\mathbf {X})\right),
\]

并且若 \(p < q\) ，则 \(\alpha\) 不是 \(f_1(\mathbf{X}) + (\mathbf{X} - \alpha)^{q - p} g_1(\mathbf{X})\) 的根；这证明了 (i)。另一方面，我们有 \(f(\mathbf{X})g(\mathbf{X}) = (\mathbf{X} - \alpha)^{p + q} f_1(\mathbf{X})g_1(\mathbf{X})\) ，且 \(f_1(\alpha)g_1(\alpha) \neq 0\) （当 \(\mathbf{A}\) 是整环时）；这证明了 (ii)。

命题5. --- 设 A 是整环。令 f 为 A{[}X{]} 的非零元素，且 \(\alpha_{1},\ldots,\alpha_{p}\) 为 f 在 A 中的两两不同的根，其阶数依次为 \(k_{1},\ldots,k_{p}\) 。我们有

\[
f (\mathrm{X}) = \left(\mathrm{X} - \alpha_ {1}\right) ^ {k _ {1}} \left(\mathrm{X} - \alpha_ {2}\right) ^ {k _ {2}} \dots \left(\mathrm{X} - \alpha_ {p}\right) ^ {k _ {p}} g (\mathrm{X})
\]

，其中 \(g \in \mathbf{A}[\mathbf{X}]\) ，且 \(\alpha_1, \ldots, \alpha_p\) 不是 \(g\) 的根。

我们对 \(p\) 作归纳，由定义 1，该命题对 \(p = 1\) 是显然的。现在设 \(f(\mathbf{X}) = g_1(\mathbf{X})g_2(\mathbf{X})\) ，其中

\[
g _ {1} (\mathrm{X}) = \left(\mathrm{X} - \alpha_ {1}\right) ^ {k _ {1}} \dots \left(\mathrm{X} - \alpha_ {p - 1}\right) ^ {k _ {p - 1}}, \quad g _ {2} (\mathrm{X}) \in \mathrm{A} [ \mathrm{X} ].
\]

由于 A 是整环且 \(\alpha_{p}\) 不同于 \(\alpha_{1},\ldots,\alpha_{p-1}\) ，因此 \(\alpha_{p}\) 不是 \(g_{1}(X)\) 的根，从而 \(\alpha_{p}\) 是 \(k_{p}\) 阶根（相对于 \(g_{2}(X)\) ；命题 4，(ii)）。由此可得 \(g_{2}(X)\) 能被 \((\mathbf{X}-\alpha_{p})^{k_{p}}\) 整除，因此

\[
f (\mathbf {X}) = \left(\mathbf {X} - \alpha_ {1}\right) ^ {k _ {1}} \dots \left(\mathbf {X} - \alpha_ {p}\right) ^ {k _ {p}} g (\mathbf {X})
\]

，其中 \(g(\mathbf{X}) \in \mathbf{A}[\mathbf{X}]\) 。显然 \(\alpha_1, \ldots, \alpha_p\) 不是 \(g\) 的根。

定理 1. --- 设 A 是一个整环。给定 A{[}X{]} 的一个非零元素 f ，其次数为 n，则 f 在 A 中所有根的阶数之和 ≤ n。这直接由命题 5 得出。

推论。--- 假设 A 是整环，且设 \(f, g \in A[X]\) ，次数 \(\leqslant n\) 。如果存在 A 的 \(n+1\) 个两两不同的元素 \(x_{1}, \ldots, x_{n+1}\) ，使得 \(f(x_{i}) = g(x_{i})\) （对 \(1 \leqslant i \leqslant n+1\) ），则 \(f = g\) 。

只需将定理1应用于 f - g 即可。

命题 6（拉格朗日插值公式）。--- 设 K 为交换域， \(\alpha_{1}, \alpha_{2}, \ldots, \alpha_{n}\) 是 K 的不同元素， \(\beta_{1}, \beta_{2}, \ldots, \beta_{n}\) 是 K 的任意元素。对 \(i = 1, 2, \ldots, n\) ，我们令

\[
f _ {i} (\mathbf {X}) = \prod_ {j \in U (i)} \left(\mathbf {X} - \alpha_ {j}\right) / \left(\alpha_ {i} - \alpha_ {j}\right),
\]

，其中 \(\mathrm{U}(i)\) 是满足 \(j \neq i\) 且 \(1 \leqslant j \leqslant n\) 的整数 j 的集合。则 \(\beta_{1}f_{1} + \cdots + \beta_{n}f_{n}\) 是 K{[}X{]} 中唯一的元素 f ，满足 \(\deg f < n\) 且 \(f(\alpha_{i}) = \beta_{i}\) （对 \(1 \leqslant i \leqslant n\) ）。

\(f\) 的唯一性由定理 1 的推论可得。设 \(f = \beta_{1}f_{1} + \cdots + \beta_{n}f_{n}\) ，则由于 \(f_{i}\) 的次数为 \(n - 1\) ，我们有 \(\deg f < n\) 。另一方面， \(f_{i}(\alpha_{j}) = 0\) （对 \(j \neq i\) ）且 \(f_{i}(\alpha_{i}) = 1\) ，因此 \(f(\alpha_{i}) = \beta_{i}\) （对 \(1 \leqslant i \leqslant n\) ）。

推论。--- 设 A 是一个整环。令 \(f \in A[X]\) ，次数 \textless{} n ，并设 K 是 A 的一个子环且是一个域。若存在 A 的 n 个不同的元素 \(\alpha_{1}, \ldots, \alpha_{n}\) ，使得 \(\alpha_{i} \in K\) 且 \(f(\alpha_{i}) \in K\) （对 \(i = 1, \ldots, n\) ），则 \(f \in K[X]\) 。

